\documentclass[aspectratio=169, 11pt]{beamer}
\usepackage{fontspec}
\setmainfont{TeX Gyre Termes}
\setsansfont{TeX Gyre Heros}
\setmonofont{Latin Modern Mono}
\usepackage[spanish,es-nodecimaldot]{babel}
\usepackage{amsmath, amssymb, amsfonts, booktabs, graphicx}
\usepackage{tikz}
\usetikzlibrary{shapes.geometric, arrows.meta, positioning, calc, fit, shadows}
\usetheme{Madrid}
\usecolortheme{beaver}
\definecolor{unalRojo}{RGB}{148, 36, 36}
\definecolor{verdeOK}{RGB}{34, 139, 34}
\definecolor{naranjaAlerta}{RGB}{230, 126, 34}
\definecolor{rojoAlerta}{RGB}{192, 57, 43}
\definecolor{azulReactor}{RGB}{41, 128, 185}
\definecolor{verdeEntrada}{RGB}{39, 174, 96}
\setbeamercolor{structure}{fg=unalRojo}
\setbeamercolor{block title}{bg=unalRojo!15, fg=unalRojo}
\setbeamercolor{block body}{bg=unalRojo!5}
\setbeamercolor{alerted text}{fg=rojoAlerta}
\newcommand{\BSF}{\textit{Hermetia illucens}}
\newcommand{\MAPE}{\text{MAPE}}
\newcommand{\DEB}{\text{DEB}}
\newcommand{\NDIR}{\text{NDIR}}
\newcommand{\MRV}{\text{MRV}}
\setbeamertemplate{navigation symbols}{}
\setbeamertemplate{footline}[frame number]

\begin{document}

% ---------------------------------------------------------------------------
% Slide 1: Preguntas consolidadas
% ---------------------------------------------------------------------------
\begin{frame}{Modelo, alcances y potenciaci\'{o}n h\'{i}brida}
    \begin{center}
        \begin{tikzpicture}
            \node[rectangle, rounded corners=5pt, draw=unalRojo, fill=unalRojo!5, 
                  text width=14cm, align=justify, inner sep=10pt, font=\small] at (0,0) {
                \textbf{P5.} Describa c\'{o}mo un modelo mecanicista basado en EDO, fundamentado en la bioenerg\'{e}tica de \BSF, puede estimar la tasa de producci\'{o}n de CO$_2$. \textbf{P8.} ?`Cu\'{a}l es la naturaleza y caracter\'{i}sticas de un modelo mecanicista y c\'{o}mo se compara con un modelo h\'{i}brido? \textbf{P3.} ?`De qu\'{e} manera la propuesta corrige o potencia el modelo de Eriksen (2022) mediante el enfoque h\'{i}brido, y por qu\'{e} CO$_2$ y CH$_4$?
            };
        \end{tikzpicture}
    \end{center}
\end{frame}

% ---------------------------------------------------------------------------
% Slide 2: Por qu\'{e} Eriksen 2022
% ---------------------------------------------------------------------------
\begin{frame}{?`Por qu\'{e} el modelo de Eriksen (2022)?}
    \begin{columns}[T]
        \begin{column}{0.55\textwidth}
            \begin{itemize}
                \item Primera integraci\'{o}n de \textbf{cin\'{e}tica logística} con \textbf{principios DEB} (Differential Energy Budget) para \BSF.
                \item Compartimenta la larva en:\\[4pt]
                \textbf{B}: biomasa estructural\\
                \textbf{L}: l\'{i}pidos de almacenamiento\\
                \textbf{A}: pool de asimilados (steady-state).
                \item Predice simult\'{a}neamente crecimiento, acumulaci\'{o}n de l\'{i}pidos y \textbf{producci\'{o}n de CO$_2$}.
            \end{itemize}
        \end{column}
        \begin{column}{0.42\textwidth}
            \begin{tikzpicture}[scale=0.75, transform shape,
                box/.style={rectangle, rounded corners=4pt, draw=unalRojo, fill=unalRojo!8, minimum width=1.6cm, minimum height=0.9cm, align=center, font=\footnotesize}]
                \node[box] (a) at (0,0) {Asimilados A};
                \node[box] (b) at (-1.8,-1.8) {Estructural B};
                \node[box] (l) at (1.8,-1.8) {L\'{i}pidos L};
                \node[box, fill=rojoAlerta!10, draw=rojoAlerta] (co2) at (0,-3.2) {CO$_2$};

                \draw[-{Stealth}, thick] (a) -- node[left, font=\tiny] {s\'{i}ntesis} (b);
                \draw[-{Stealth}, thick] (a) -- node[right, font=\tiny] {almacenamiento} (l);
                \draw[-{Stealth}, thick] (a) -- node[right, font=\tiny] {mantenimiento} (co2);
                \draw[-{Stealth}, thick] (b) -- node[left, font=\tiny] {costo} (co2);
                \draw[-{Stealth}, thick] (l) -- node[right, font=\tiny] {costo} (co2);
                \draw[-{Stealth}, thick, dashed] (l) to[out=120, in=60] node[above, font=\tiny] {reciclaje} (a);
            \end{tikzpicture}
        \end{column}
    \end{columns}

    \vfill
    \begin{block}{Ecuaci\'{o}n de CO$_2$ (Eriksen)}
        \centering
        $r_{\text{CO}_2} = \underbrace{m \cdot B}_{\text{mantenimiento}} + \underbrace{\frac{1-Y_B}{Y_B} \cdot r_B}_{\text{s\'{i}ntesis } B} + \underbrace{\frac{1-Y_L}{Y_L} \cdot r_L}_{\text{s\'{i}ntesis } L}$
    \end{block}
\end{frame}

% ---------------------------------------------------------------------------
% Slide 3: Supuestos del modelo
% ---------------------------------------------------------------------------
\begin{frame}{Supuestos del modelo mecanicista}
    \begin{columns}[T]
        \begin{column}{0.48\textwidth}
            \begin{block}{Supuestos biof\'{i}sicos}
                \begin{itemize}
                    \item \textbf{Una sola larva}: individuo promedio, no poblaci\'{o}n.
                    \item \textbf{6 instares}: transiciones modeladas por peso umbral ($\delta_I = 5$).
                    \item Asimilados en \textbf{steady-state}: pool A peque\~{n}o y constante.
                    \item \textbf{Homogeneidad espacial}: par\'{a}metros concentrados, sin gradientes 3D.
                \end{itemize}
            \end{block}
        \end{column}
        \begin{column}{0.48\textwidth}
            \begin{block}{Supuestos cin\'{e}ticos}
                \begin{itemize}
                    \item Asimilaci\'{o}n limitada por log\'{i}stica de peso (instar 6).
                    \item Crecimiento estructural con segunda log\'{i}stica ($\beta$) que genera surplus de l\'{i}pidos.
                    \item Mantenimiento proporcional a \textbf{B}, no a L.
                    \item Reciclaje de l\'{i}pidos cuando A $<$ mantenimiento.
                \end{itemize}
            \end{block}
        \end{column}
    \end{columns}
\end{frame}

% ---------------------------------------------------------------------------
% Slide 4: De una larva al reactor
% ---------------------------------------------------------------------------
\begin{frame}{De una larva al reactor: ?`por qu\'{e} es \'{u}til?}
    \begin{columns}[T]
        \begin{column}{0.55\textwidth}
            \begin{alertblock}{El modelo es individual, pero...}
                \begin{itemize}
                    \item La fisiolog\'{i}a individual escala linealmente con la densidad larval en condiciones no limitantes.
                    \item Proporciona una \textbf{l\'{i}nea base termodin\'{a}micamente consistente} contra la cual comparar.
                    \item Permite calcular coeficientes de rendimiento estequiom\'{e}tricos ($Y_{\text{CO}_2}$) v\'{a}lidos por gramo de biomasa.
                    \item Es la \textbf{restricci\'{o}n f\'{i}sica} que todo h\'{i}brido debe respetar.
                \end{itemize}
            \end{alertblock}
        \end{column}
        \begin{column}{0.42\textwidth}
            \begin{center}
                \begin{tikzpicture}[scale=0.8, transform shape,
                    box/.style={rectangle, rounded corners=4pt, draw=azulReactor, fill=azulReactor!12, minimum width=2.4cm, minimum height=1cm, align=center, font=\small}]
                    \node[box] (ind) at (0,1.5) {Larva individual\\Eriksen 2022};
                    \node[box] (esc) at (0,0) {Escalamiento\\$\times$ densidad};
                    \node[box] (rea) at (0,-1.5) {Reactor real\\+ correcci\'{o}n h\'{i}brida};

                    \draw[-{Stealth}, thick, azulReactor] (ind) -- (esc);
                    \draw[-{Stealth}, thick, azulReactor] (esc) -- (rea);
                \end{tikzpicture}
            \end{center}
        \end{column}
    \end{columns}
\end{frame}

% ---------------------------------------------------------------------------
% Slide 5: Limitaciones estructurales
% ---------------------------------------------------------------------------
\begin{frame}{Limitaciones estructurales del mecanicista}
    \begin{columns}[T]
        \begin{column}{0.48\textwidth}
            \begin{alertblock}{Ceguera ante la heterogeneidad}
                \begin{itemize}
                    \item Agregaciones larvarias = fluidos activos (gradientes T, O$_2$).
                    \item Atribuye desviaciones a ``error param\'{e}trico'' cuando son arquitect\'{o}nicas.
                    \item No predice CH$_4$: carece de compartimento para anaerobiosis.
                \end{itemize}
            \end{alertblock}
        \end{column}
        \begin{column}{0.48\textwidth}
            \begin{alertblock}{Rigidez y validaci\'{o}n}
                \begin{itemize}
                    \item Par\'{a}metros ($\alpha$, $\beta$, $a_{\text{max}}$) espec\'{i}ficos por sustrato.
                    \item Validaci\'{o}n destructiva: muestrear biomasa altera el reactor.
                    \item Ruido correlacionado (condensaci\'{o}n, obstrucci\'{o}n) se confunde con variabilidad biol\'{o}gica.
                \end{itemize}
            \end{alertblock}
        \end{column}
    \end{columns}

    \vfill
    \begin{center}
        \textbf{El valor no es la predicci\'{o}n puntual, sino la l\'{i}nea base f\'{i}sicamente consistente.}
    \end{center}
\end{frame}

% ---------------------------------------------------------------------------
% Slide 6: Potenciaci\'{o}n h\'{i}brida
% ---------------------------------------------------------------------------
\begin{frame}{Potenciaci\'{o}n h\'{i}brida: dos frentes complementarios}
    \begin{center}
        \begin{tikzpicture}[scale=0.82, transform shape,
            box/.style={rectangle, rounded corners=5pt, draw=unalRojo, fill=unalRojo!5, minimum width=3.2cm, minimum height=1.3cm, align=center, font=\small},
            arrow/.style={-{Stealth}, line width=2pt, unalRojo}]

            \node[box] (sen) at (-4.5,0) {\textbf{Sensor virtual}\\Red densa (FFN)};
            \node[box] (pinn) at (0,0) {\textbf{PINN}\\Eriksen + datos};
            \node[box] (out) at (4.5,0) {\textbf{Predicci\'{o}n}\\F\'{i}sicamente consistente};

            \draw[arrow] (sen) -- node[above, font=\footnotesize] {$B(t)$ estimada} (pinn);
            \draw[arrow] (pinn) -- (out);
        \end{tikzpicture}
    \end{center}

    \vfill
    \begin{columns}[T]
        \begin{column}{0.48\textwidth}
            \begin{block}{Sensor virtual}
                \begin{itemize}
                    \item Entradas: densidad inicial, T, humedad.
                    \item Salida: $B(t)$ sin muestreo destructivo.
                    \item Actualiza $dB/dt$ en cada paso de tiempo.
                \end{itemize}
            \end{block}
        \end{column}
        \begin{column}{0.48\textwidth}
            \begin{block}{PINN}
                \begin{itemize}
                    \item P\'{e}rdida dual: error de datos + residual EDO.
                    \item Conservaci\'{o}n de masa como \textbf{restricci\'{o}n dura}.
                    \item Captura no linealidades (agregaciones, anaerobiosis).
                \end{itemize}
            \end{block}
        \end{column}
    \end{columns}
\end{frame}

% ---------------------------------------------------------------------------
% Slide 7: Justificaci\'{o}n CO2/CH4 y cierre
% ---------------------------------------------------------------------------
\begin{frame}{Justificaci\'{o}n del alcance: CO$_2$ y CH$_4$}
    \begin{columns}[T]
        \begin{column}{0.48\textwidth}
            \begin{block}{CO$_2$: se\~{n}al fisiol\'{o}gica}
                \begin{itemize}
                    \item Gas primario de la respiraci\'{o}n aerobia.
                    \item Directamente vinculado a $B(t)$ por el coeficiente $Y_{\text{CO}_2}$ de Eriksen.
                    \item Condici\'{o}n necesaria para validar el estado metab\'{o}lico y calibrar el h\'{i}brido.
                \end{itemize}
            \end{block}
        \end{column}
        \begin{column}{0.48\textwidth}
            \begin{alertblock}{CH$_4$: se\~{n}al de fallo}
                \begin{itemize}
                    \item No es metabolismo larval: es metanog\'{e}nesis microbiana en microambientes an\'{o}xicos.
                    \item Indicador operativo de anaerobiosis (exceso humedad, falla aireaci\'{o}n).
                    \item PGW 27,2$\times$ CO$_2$ (IPCC AR6): prioridad clim\'{a}tica.
                \end{itemize}
            \end{alertblock}
        \end{column}
    \end{columns}

    \vfill
    \begin{center}
        \textbf{N$_2$O queda fuera del alcance cuantitativo:} requiere sensores electroqu\'{i}micos espec\'{i}ficos y cadenas de calibraci\'{o}n independientes, fuera del dominio de validaci\'{o}n de esta fase.
    \end{center}
\end{frame}

\end{document}
